% Generated draft: requires mathematical and source review.
\documentclass[11pt]{article}
\usepackage[margin=25mm]{geometry}
\usepackage{fontspec}
\setmainfont{Latin Modern Roman}
\usepackage{amsmath,amssymb,mathtools,unicode-math}
\setmathfont{Latin Modern Math}
\usepackage{enumitem,needspace,xurl,hyperref}
\setlength{\parindent}{0pt}
\setlength{\parskip}{5pt}
\setlength{\emergencystretch}{3em}
\newcommand{\N}{\mathbb{N}}
\newcommand{\Z}{\mathbb{Z}}
\newcommand{\Q}{\mathbb{Q}}
\newcommand{\R}{\mathbb{R}}
\newcommand{\C}{\mathbb{C}}
\newcommand{\F}{\mathbb{F}}
\newcommand{\A}{\mathbb{A}}
\newcommand{\PP}{\mathbb{P}}
\newcommand{\E}{\mathbb{E}}
\newcommand{\eps}{\varepsilon}
\newcommand{\dd}{\,\mathrm{d}}
\begin{document}
\section*{UnsolvedMath ID: 701134}
\textbf{Source problem code:} catalogue entry
\subsection{Definitions and mathematical statement}
Let \(K\) be a field, let \(\Phi\) be the root system occurring in the source, and let \(G_\Phi(K)\) denote the associated Chevalley group. For \(r\in\Phi\), write \(x_r: (K,+)\to G_\Phi(K)\) for its root subgroup parametrization and, for \(B\subseteq K\), put \(x_r(B)=\{x_r(b):b\in B\}\).

Let \(\mathfrak A=\{\mathfrak A_r\mid r\in\Phi\}\) be an elementary carpet of additive subgroups of \(K\), in the sense intended by the source, and define
\[
E(\mathfrak A)=\left\langle x_r(\mathfrak A_r)\mid r\in\Phi\right\rangle\leq G_\Phi(K).
\]
The carpet is called \emph{admissible} if
\[
E(\mathfrak A)\cap x_r(K)=x_r(\mathfrak A_r)
\qquad\text{for every }r\in\Phi.
\]
The supplied record does not define the additional conditions expressed by the term \emph{elementary carpet}; consequently, the admissible class of input families cannot be fully specified from this record alone.

For each \(r\in\Phi\), put
\[
G_r=\langle x_r(K),x_{-r}(K)\rangle,
\qquad
E_r(\mathfrak A)=\langle x_r(\mathfrak A_r),x_{-r}(\mathfrak A_{-r})\rangle\leq G_r.
\]
Interpreting the rank-one subcarpet inside this ambient rank-one subgroup, its admissibility means
\[
E_r(\mathfrak A)\cap x_s(K)=x_s(\mathfrak A_s)
\qquad\text{for each }s\in\{r,-r\}.
\]
The question is whether, for every field \(K\), every root system \(\Phi\) allowed by the source, and every elementary carpet \(\mathfrak A\) of type \(\Phi\), one has
\[
\begin{aligned}
&\bigl[\,E(\mathfrak A)\cap x_r(K)=x_r(\mathfrak A_r)
       \text{ for every }r\in\Phi\,\bigr]\\
&\qquad\Longleftrightarrow
\bigl[\,E_r(\mathfrak A)\cap x_s(K)=x_s(\mathfrak A_s)
       \text{ for every }r\in\Phi\text{ and every }s\in\{r,-r\}\,\bigr].
\end{aligned}
\]
This is the proposed reduction to Lie rank one of the admissibility question numbered 7.28 in the Kourovka Notebook. That earlier question is not supplied. The precise scope of \(\Phi\), the form of the Chevalley group, and the intended rank-one realization remain source qualifications requiring confirmation.

The supplied record reports that the answer is affirmative when \(K\) is locally finite, meaning that every finitely generated subfield of \(K\) is finite, and attributes this result to V. M. Levchuk (1983). This is a supplied literature assertion, not an independently verified result here.
\subsection{Short English statement}
V. M. Levchuk asks whether an elementary carpet over an arbitrary field is admissible exactly when each of its opposite-root rank-one subcarpets is admissible. The supplied source reports an affirmative answer over locally finite fields. The definition of elementary carpet and some ambient group conventions are missing, so the statement requires human mathematical review.
\subsection{Sources}
\begin{itemize}
\item Local UnsolvedMath registry: supplied record with permanent ID \(701134\), exact source code \texttt{catalogue entry}, titled \emph{Kourovka Notebook Problem 15.46}. The registry labels the problem open; this is supplied status metadata, not an independently verified assessment.
\item \emph{The Kourovka Notebook (Unsolved Problems in Group Theory)}, problem 15.46, proposed by V. M. Levchuk in the 15th issue (2002). The supplied background also identifies the 21st issue edition (2026), \texttt{arXiv:1401.0300v48}; the collection metadata names E. I. Khukhro and V. D. Mazurov as editors.
\item Kourovka Notebook question 7.28 is referenced as the underlying admissibility question; its text and detailed citation are not supplied.
\item V. M. Levchuk, \emph{Algebra and Logic}, volume 22, no. 5 (1983), pp. 362--371. The supplied record cites this reference for the affirmative answer over locally finite fields; no article title is supplied.
\end{itemize}
\end{document}
